\documentclass[a4paper,12pt]{article}
\usepackage[T2A]{fontenc}
\usepackage[utf8]{inputenc}
\usepackage[russian]{babel}
\usepackage{amsmath,amssymb,mathtools}
\usepackage{helvet}
\renewcommand{\familydefault}{\sfdefault}
\usepackage{icomma}
\usepackage[a4paper,left=23mm,right=23mm,top=23mm,bottom=23mm,headheight=15pt]{geometry}
\usepackage{microtype}
\usepackage{tikz}
\usetikzlibrary{calc,arrows.meta,positioning,shapes.geometric}
\usepackage{tabularx,booktabs}
\usepackage{enumitem}
\usepackage{fancyhdr}
\usepackage{setspace}
\usepackage{xcolor}
\usepackage{hyperref}
\hypersetup{hidelinks,pdftitle={Чек-лист. Логарифмические функции и уравнения},pdfauthor={Лёвин Артём Александрович}}
\definecolor{accent}{HTML}{137A7A}
\definecolor{darkaccent}{HTML}{0A525A}
\definecolor{textgray}{HTML}{2F383B}
\definecolor{rulegray}{HTML}{C5D4D4}
\color{textgray}
\onehalfspacing
\setlength{\parindent}{0pt}
\setlength{\parskip}{3pt}
\setcounter{secnumdepth}{0}
\setlist[enumerate]{leftmargin=1.9em,itemsep=5pt,topsep=6pt}
\pagestyle{fancy}
\fancyhf{}
\fancyhead[L]{\small\color{darkaccent}Марк Гукин\quad\textbar\quad 21.10.2026}
\fancyhead[R]{\small\color{textgray}Логарифмы}
\fancyfoot[C]{\small\color{textgray}\thepage}
\renewcommand{\headrulewidth}{0.35pt}
\renewcommand{\footrulewidth}{0pt}
\newcommand{\sect}[1]{\par\vspace{3.5mm}{\Large\bfseries\color{darkaccent}#1}\par\vspace{1mm}\color{textgray}}
\newcommand{\subj}[1]{\par\vspace{1.5mm}{\large\bfseries\color{darkaccent}#1}\par\color{textgray}}
\newcommand{\key}[1]{\textbf{\color{darkaccent}#1}}
\newcommand{\checkitem}[1]{\par\noindent$\square$\enspace #1\par}
\begin{document}

% Титульный лист
\thispagestyle{empty}
\begin{center}
\vspace*{35mm}
{\small\bfseries\color{darkaccent}ИНДИВИДУАЛЬНОЕ УЧЕБНОЕ ПОСОБИЕ}\par
\vspace{18mm}
{\fontsize{34}{40}\selectfont\bfseries\color{darkaccent}Чек-лист}\par
\vspace{9mm}
{\Large\bfseries Логарифмические функции}\par
{\Large\bfseries и уравнения}\par
\vspace{6mm}
{\large Метод двух точек, степени и ОДЗ}\par
\vspace{21mm}
{\large\bfseries Марк Гукин}\par
\vspace{8mm}
{\normalsize 21 октября 2026 года}\par
\vfill
{\normalsize Лёвин Артём Александрович}\par
{\normalsize\bfseries\color{darkaccent}эксклюзивно для Марка Гукина}\par
\vspace{17mm}
\begin{tikzpicture}[x=1cm,y=1cm]
  \draw[darkaccent,opacity=.42,line width=.9pt]
  (-7.4,0) .. controls (-5.0,.30) and (-3.6,-.38) .. (-.3,.13)
  .. controls (2.4,.47) and (4.5,-.31) .. (7.4,.05);
\end{tikzpicture}
\vspace{8mm}
\end{center}
\clearpage

\sect{01\quad Логарифм и обратный переход}
\checkitem{Понимаю определение логарифма и всегда проверяю условия его существования.}
\[
\boxed{\log_a b=c\ \Longleftrightarrow\ b=a^c,
\qquad a>0,\quad a\ne1,\quad b>0.}
\]
\key{«Метод циклопа»} --- запоминающаяся схема: \emph{аргумент логарифма равен основанию в степени, стоящей справа}. Это обычный переход от логарифмической записи к показательной.
\[
\log_2\frac{80}{t-20}=3
\quad\Longleftrightarrow\quad
\frac{80}{t-20}=2^3=8
\quad\Longrightarrow\quad \boxed{t=30}.
\]
Здесь \(t>20\) --- условие существования исходного логарифма. При \(t=30\) оно выполняется.

\sect{02\quad Подготовка чисел и вычисления}
\checkitem{Записываю десятичные дроби и большие числа через степени десяти, чтобы упростить подстановку.}
\[
0{,}7=7\cdot10^{-1},\qquad
0{,}4=4\cdot10^{-1},\qquad
4200=42\cdot10^2.
\]
\key{Пример.} Найдите \(t\), если
\[
84=\frac{0{,}7\cdot4200\cdot0{,}4}{42}
\log_2\frac{100-20}{t-20}.
\]
\[
\begin{aligned}
\frac{0{,}7\cdot4200\cdot0{,}4}{42}
&=\frac{(7\cdot10^{-1})(42\cdot10^2)(4\cdot10^{-1})}{42}=28,\\
84&=28\log_2\frac{80}{t-20},\\
\log_2\frac{80}{t-20}&=3,\\
\frac{80}{t-20}&=8,\qquad \boxed{t=30}.
\end{aligned}
\]
\key{Проверка:} \(t-20=10>0\); \(\log_2(80/10)=\log_2 8=3\); \(28\cdot3=84\).

\sect{03\quad Свойства, которые понадобятся}
\checkitem{Узнаю степени и рационально применяю свойства логарифмов.}
\[
\log_a1=0,\quad \log_a a=1,\quad
\log_a (u^n)=n\log_a u\quad (u>0).
\]
\[
\log_{a^p}u=\frac1p\log_a u\quad(p\ne0),\qquad
\log_a u+\log_a v=\log_a(uv)\quad(u,v>0).
\]
Например, \(\log_{\sqrt5}25=\dfrac{2}{1/2}=4\), а \(1=\log_5 5\). В этих формулах все основания положительны и отличны от единицы.
\clearpage

\sect{04\quad Восстанавливаем функцию по графику}
\checkitem{Выбираю две удобные узловые точки для двух неизвестных параметров.}
\key{Узловая точка} --- точка с целочисленными координатами на изображённом графике. При двух неизвестных параметрах достаточно двух подходящих точек; особенно удобны координаты \(x=1\), \(y=0\), а также небольшие целые числа.

\subj{Функция \(f(x)=b+\log_a x\)}
Пусть на графике лежат точки \(A(1;-3)\) и \(B(4;-1)\). Тогда
\[
\begin{cases}
-3=b+\log_a1,\\
-1=b+\log_a4,
\end{cases}
\quad\Longrightarrow\quad
b=-3,\quad \log_a4=2,\quad a^2=4.
\]
Так как \(a>0\), \(a\ne1\), получаем \(a=2\) и
\[
f(x)=-3+\log_2x,\qquad
\boxed{f(32)=-3+5=2}.
\]
\key{Область определения:} \(x>0\).

\subj{Функция \(f(x)=\log_a(x+b)\)}
Пусть на графике лежат \(A(-4;0)\) и \(B(-1;2)\). Подставим координаты:
\[
\begin{cases}
0=\log_a(-4+b),\\
2=\log_a(-1+b).
\end{cases}
\quad\Longleftrightarrow\quad
\begin{cases}
-4+b=a^0=1,\\
-1+b=a^2.
\end{cases}
\]
Из первого равенства \(b=5\), из второго \(a^2=4\). С учётом \(a>0\) имеем \(a=2\). Следовательно,
\[
f(x)=\log_2(x+5),\qquad
\boxed{f(11)=\log_2 16=4}.
\]
\key{Область определения:} \(x+b>0\); в примере \(x>-5\).

\sect{05\quad Если основание получилось дробным или корнем}
\checkitem{Переписываю основание в степенном виде и правильно выношу показатель.}
\[
\log_{1/\sqrt3}27
=\log_{3^{-1/2}}3^3
=\frac{3}{-1/2}=\boxed{-6}.
\]
\key{Следите за направлением:} показатель в основании переходит в \emph{знаменатель}. Основание логарифма может быть между нулём и единицей; отрицательным или равным единице оно быть не может.
\clearpage

\sect{06\quad Решаем логарифмические уравнения}
\checkitem{Перед решением отмечаю ОДЗ каждого логарифма, после решения проверяю найденные значения.}
\[
\boxed{\log_a U(x)=c\quad\Longleftrightarrow\quad U(x)=a^c}
\]
При этом \(a>0\), \(a\ne1\), \(U(x)>0\). В случае \(\log_aU(x)=\log_aV(x)\) можно перейти к \(U(x)=V(x)\), сохраняя \(U(x)>0\) и \(V(x)>0\).

\subj{Пример 1. Дробное основание}
\[
\log_{1/7}(7-x)=-2
\quad\Longleftrightarrow\quad
7-x=(7^{-1})^{-2}=49
\quad\Longrightarrow\quad \boxed{x=-42}.
\]
\key{Проверка ОДЗ:} \(7-(-42)=49>0\). Найденный корень подходит.

\subj{Пример 2. Справа логарифм и число}
\[
\log_5(7-x)=\log_5(3-x)+1.
\]
\key{ОДЗ:} \(7-x>0\) и \(3-x>0\), то есть \(x<3\). Запишем \(1=\log_5 5\):
\[
\begin{aligned}
\log_5(7-x)
&=\log_5(3-x)+\log_55\\
&=\log_5\bigl(5(3-x)\bigr),\\
7-x&=5(3-x),\\
4x&=8,\qquad \boxed{x=2}.
\end{aligned}
\]
\key{Проверка ОДЗ:} \(7-2=5>0\), \(3-2=1>0\). Оба логарифма определены.

\subj{Самопроверка перед записью ответа}
\checkitem{Для каждого логарифма проверено положительное основание, отличное от единицы.}
\checkitem{Для каждого логарифма проверена положительность аргумента.}
\checkitem{Все посторонние значения отброшены; вычисления проверены подстановкой.}
\clearpage

% Блок-схема алгоритма на отдельной странице A4.
\thispagestyle{plain}
\begin{center}
{\Large\bfseries\color{darkaccent}Алгоритм: логарифмическое уравнение}\par
\vspace{11mm}
\begin{tikzpicture}[
  >=Latex,
  node distance=12mm,
  terminator/.style={ellipse,draw=darkaccent,thick,align=center,minimum width=45mm,minimum height=11mm,font=\small},
  action/.style={rectangle,rounded corners=2pt,draw=darkaccent,thick,align=center,text width=88mm,minimum height=14mm,inner sep=4pt,font=\small},
  decision/.style={diamond,aspect=2.6,draw=darkaccent,thick,align=center,text width=54mm,inner sep=3pt,font=\small},
  arrow/.style={->,line width=.85pt,draw=darkaccent}
]
\node[terminator] (s) {Начало};
\node[action,below=of s] (dom) {Запишите ОДЗ всех логарифмов};
\node[decision,below=13mm of dom] (q) {После преобразований\\получились два логарифма\\с одним основанием?};
\node[action,below left=20mm and -5mm of q,text width=56mm] (left) {Приравняйте аргументы этих логарифмов};
\node[action,below right=20mm and -5mm of q,text width=56mm] (right) {Если \(\log_a U=c\), перейдите к \(U=a^c\)};
\node[action,below=40mm of q] (solve) {Решите полученное алгебраическое уравнение};
\node[action,below=of solve] (check) {Подставьте все кандидаты в ОДЗ и проверьте исходное равенство};
\node[terminator,below=of check] (end) {Запишите допустимые корни};
\draw[arrow] (s)--(dom);
\draw[arrow] (dom)--(q);
\draw[arrow] (q.west) -- node[above,font=\scriptsize]{да} ++(-10mm,0) |- ([yshift=9mm]left.north) -- (left.north);
\draw[arrow] (q.east) -- node[above,font=\scriptsize]{нет} ++(10mm,0) |- ([yshift=9mm]right.north) -- (right.north);
\draw[arrow] (left.south) |- (solve.west);
\draw[arrow] (right.south) |- (solve.east);
\draw[arrow] (solve)--(check);
\draw[arrow] (check)--(end);
\end{tikzpicture}
\end{center}
\clearpage

\sect{Тренировочный блок\quad Путь А}
\textit{10 заданий. Выполняйте по порядку; в уравнениях записывайте ОДЗ и проверяйте корни.}
\vspace{1mm}
\begin{enumerate}[itemsep=7pt]
\item Представьте \(0{,}6\), \(0{,}8\) и \(3600\) через степени десяти.\\
Вычислите \(\displaystyle\frac{0{,}6\cdot3600\cdot0{,}8}{36}\).
\item Решите уравнение \(\displaystyle\log_2\frac{72}{t-12}=3\).
\item Найдите \(t\), если
\(\displaystyle 48=\frac{0{,}6\cdot3600\cdot0{,}8}{36}\log_2\frac{88-8}{t-8}\).
\item График функции \(f(x)=b+\log_a x\) проходит через точки \((1;-2)\) и \((9;0)\). Найдите \(f(27)\).
\item График функции \(f(x)=\log_a(x+b)\) проходит через точки \((-2;0)\) и \((2;2)\). Найдите \(f(22)\).
\item Решите уравнение \(\displaystyle\log_{1/5}(9-x)=-2\).
\item Решите уравнение \(\displaystyle\log_4(2x+5)=2\).
\item Решите уравнение \(\displaystyle\log_3(13-x)=\log_3(4-x)+1\).
\item Решите уравнение \(\displaystyle\log_2(x-6)=\log_2(2x-10)+1\). Особое внимание уделите ОДЗ.
\item График функции \(f(x)=b+\log_a x\) проходит через точки \((1;3)\) и \(\bigl(\frac14;5\bigr)\). Найдите \(f(8)\).
\end{enumerate}
\clearpage

\sect{Ответы к тренировочному блоку}
\renewcommand{\arraystretch}{1.55}
\begin{tabularx}{\linewidth}{@{}c X@{}}
\toprule
\textbf{№} & \textbf{Ответ}\\
\midrule
1 & \(0{,}6=6\cdot10^{-1}\), \(0{,}8=8\cdot10^{-1}\), \(3600=36\cdot10^2\); значение \(48\).\\
2 & \(t=21\).\\
3 & \(t=48\).\\
4 & \(f(27)=1\), \(a=3\), \(b=-2\).\\
5 & \(f(22)=4\), \(a=\sqrt5\), \(b=3\).\\
6 & \(x=-16\).\\
7 & \(x=\dfrac{11}{2}\).\\
8 & \(x=-\dfrac12\).\\
9 & Корней нет: кандидат \(x=\dfrac{14}{3}\) нарушает ОДЗ \(x>6\).\\
10 & \(f(8)=0\), \(a=\dfrac12\), \(b=3\).\\
\bottomrule
\end{tabularx}
\par\vspace{11mm}
\noindent\key{Короткая памятка.} Если равенство после преобразований привело к числу, нарушающему хотя бы одно исходное ограничение, это число не является корнем уравнения.
\end{document}
